\documentclass[11pt,a4paper]{article}
\usepackage[utf8]{inputenc}
\usepackage[T1]{fontenc}
\usepackage{amsmath,amssymb}
\usepackage{booktabs}
\usepackage{multirow}
\usepackage{graphicx}
\usepackage{hyperref}
\usepackage[margin=2.5cm]{geometry}
\usepackage{xcolor}
\usepackage{tabularx}

\title{Seed Variance Dominates Error Injection Signal\\in Low-Data Embodied Planning SFT}

\author{Anonymous}

\date{}

\begin{document}
\maketitle

\begin{abstract}
Taxonomy-guided error injection---exposing LLM agents to structured failure modes during supervised fine-tuning (SFT)---has shown large gains in tool-use domains.
We test whether this approach transfers to embodied planning by injecting errors from a 4-operator taxonomy into ALFWorld SFT trajectories.
Initial single-seed results suggested a variance gradient, a threshold-plateau dose-response, and a constant ${\sim}6.4$pp taxonomy penalty---patterns that would imply task-structural boundary conditions for error injection.
Multi-seed evaluation reveals that \emph{all three findings are artifacts of seed selection}.
Across three seeds the gap between taxonomy-guided injection (TaxEI) and vanilla SFT is $-8.6$pp but non-significant ($p{=}0.32$); per-operator advantages observed at seed\,=\,42 vanish; and scale ablation to 25K--50K trajectories closes the remaining gap to ${\leq}2$pp.
A second environment (WebShop) confirms the null.
Training dynamics analysis shows the model absorbs injected and clean samples at identical rates; qualitative failure analysis finds no injection-correlated error patterns ($\chi^2$ $p{=}0.91$).
We draw two methodological conclusions: (1) cross-seed variance in small-sample agent SFT (${\leq}10$K trajectories) is so large that data augmentation evaluations require ${\geq}5$ seeds to be informative, and (2) below this scale, error injection signal is indistinguishable from training noise regardless of injection strategy.
\end{abstract}

%------------------------------------------------------------------
\section{Introduction}
%------------------------------------------------------------------

Data augmentation through error injection is an intuitive training strategy for LLM agents: expose the model to failure modes so it learns to avoid or recover from them.
PALADIN \cite{paladin2025} demonstrated this for tool-use agents on ToolBench, using a 7-type failure taxonomy to generate 50K+ recovery-annotated trajectories and achieving a +57pp Recovery Rate gain (though ${\sim}30$pp of this is attributable to inference-time retrieval augmentation rather than SFT alone).

A natural follow-up question is whether the error-injection \emph{training} component transfers to embodied planning, where agents execute multi-step plans in simulated physical environments.
We designed TaxEI (Taxonomy-guided Error Injection), adapting the error-injection paradigm to ALFWorld with a 4-operator taxonomy covering step-order, sub-goal, context, and observation perturbations (\S\ref{sec:method}).

Our initial evaluation at a single seed (seed\,=\,42) produced three seemingly interesting patterns: (i) a variance gradient where more structured injection produced more instability, (ii) a threshold-plateau dose-response between 10\% and 25\% injection ratio, and (iii) a remarkably constant ${\sim}6.4$pp penalty from taxonomy labels compared to random injection.
Each suggested that task structure---specifically the difference between tool-use and embodied planning---determines whether error injection helps or hurts.

\textbf{All three findings dissolved under multi-seed evaluation.}
Expanding from one to three seeds, the TaxEI--Vanilla gap becomes non-significant ($p{=}0.32$).
Per-operator advantages observed at seed\,=\,42 (e.g., C4 observation masking at +9.1pp) vanish at three-seed means.
Scale ablation to 25K and 50K trajectories closes the gap further, confirming that the original findings were low-data artifacts rather than structural boundary conditions.

This paper therefore makes two contributions:

\begin{enumerate}
\item \textbf{Evaluation methodology}: Cross-seed standard deviation in ALFWorld SFT reaches 5--8pp across methods.
  At this variance level, any data augmentation claim based on fewer than five seeds is unreliable.
  We provide concrete evidence for this via six experiments whose single-seed conclusions all reverse or dissolve under multi-seed aggregation.

\item \textbf{Scale-dependent null result}: Below ${\sim}10$K training trajectories, error injection signal---whether taxonomy-guided or random---is indistinguishable from training noise in embodied planning.
  This null persists across two environments (ALFWorld, WebShop), multiple injection strategies, and both SFT and in-context learning.
\end{enumerate}

We believe this constitutes a useful negative result for the community: it redirects effort away from low-data error injection in planning domains and provides a concrete seed-count recommendation for future data augmentation evaluations.

%------------------------------------------------------------------
\section{Method: TaxEI}\label{sec:method}
%------------------------------------------------------------------

\paragraph{Environment.}
ALFWorld is a text-based embodied planning benchmark where an agent navigates household environments to complete tasks (e.g., ``heat a mug and place it on the shelf'').
The agent interacts through natural language actions and receives textual observations.
Crucially, syntactically valid but semantically wrong actions produce valid-looking observations---errors are \emph{silent}.

\paragraph{Error taxonomy.}
We designed four perturbation operators targeting distinct failure modes in embodied planning:
\begin{itemize}
\item \textbf{P1} (step\_order\_swap): Permute the ordering of actions within a trajectory.
\item \textbf{P2} (sub\_goal\_deletion): Remove sub-goals from successful trajectories.
\item \textbf{C1} (context\_truncation): Truncate environmental context provided to the agent.
\item \textbf{C4} (observation\_masking): Mask observation feedback from the environment.
\end{itemize}

\paragraph{Data construction.}
From 501 GPT-5.2-generated successful ALFWorld trajectories, we produced 4,671 injected variants using the four operators in round-robin fashion.
Three SFT conditions were constructed:
\begin{itemize}
\item \textbf{Vanilla}: Clean trajectories only (4,358 samples).
\item \textbf{TaxEI}: Clean + taxonomy-tagged injections (9,029 samples). Each injected sample carries a label indicating which operator produced it.
\item \textbf{Random}: Clean + unlabeled injections (9,029 samples). Injected samples carry a generic error marker without taxonomy information.
\end{itemize}
TaxEI and Random are matched in data volume and injection ratio; they differ only in whether the error tag encodes the taxonomy label.

\paragraph{Training.}
All conditions use Qwen3-8B as the base model, trained for 3 epochs with identical hyperparameters.
Each condition is trained with three seeds (42, 0, 123) and evaluated on all 47 ALFWorld test tasks.

%------------------------------------------------------------------
\section{Experiments}
%------------------------------------------------------------------

We present six experiments, each designed to probe a different facet of the null result.
All quantitative claims are derived from the data; $p$-values are from two-sided paired $t$-tests unless otherwise noted.

%------------------------------------------------------------------
\subsection{Main Results: Multi-Seed Evaluation}\label{sec:main}
%------------------------------------------------------------------

Table~\ref{tab:main} presents three-seed results for all three conditions.

\begin{table}[h]
\centering
\caption{ALFWorld 3-seed clean task performance.}
\label{tab:main}
\begin{tabular}{lcccc}
\toprule
\textbf{Method} & \textbf{Mean} & \textbf{SD} & \textbf{vs.\ Vanilla} & \textbf{$p$-value} \\
\midrule
Vanilla & 65.3\% & 3.3pp & --- & --- \\
Random  & 63.1\% & 7.0pp & $-2.2$pp & NS \\
TaxEI   & 56.7\% & 8.6pp & $-8.6$pp & 0.32 \\
\bottomrule
\end{tabular}
\end{table}

The TaxEI--Vanilla gap of $-8.6$pp appears large in absolute terms but is not statistically significant ($p{=}0.32$) due to extreme cross-seed variance: TaxEI achieves 61.7\% at seeds 42 and 123 (matching Vanilla seed\,=\,42 exactly) but collapses to 46.8\% at seed\,=\,0---a 14.9pp swing within the same method.
The variance ratio TaxEI/Vanilla is 2.6$\times$, but this ratio is itself unstable at $n{=}3$ seeds.

\textbf{Key observation}: At 3 seeds with SDs of 3--9pp, \emph{no pairwise comparison achieves significance}.
A single favorable or unfavorable seed can shift the mean by ${\sim}8$pp, making any conclusion from fewer seeds unreliable.

%------------------------------------------------------------------
\subsection{Training Dynamics Analysis}\label{sec:dynamics}
%------------------------------------------------------------------

If error injection corrupts training, injected samples should exhibit systematically different loss trajectories from clean samples.
Table~\ref{tab:loss} shows per-epoch loss decomposed by sample type.

\begin{table}[h]
\centering
\caption{Training loss by sample type (3-seed mean $\pm$ SE).}
\label{tab:loss}
\begin{tabular}{llcccc}
\toprule
\textbf{Method} & \textbf{Sample Type} & \textbf{Ep.\,1} & \textbf{Ep.\,2} & \textbf{Ep.\,3} & \textbf{Seed SD} \\
\midrule
TaxEI   & Clean    & 0.58\,$\pm$\,0.03 & 0.38\,$\pm$\,0.02 & 0.29\,$\pm$\,0.02 & 0.04 \\
TaxEI   & Injected & 0.61\,$\pm$\,0.04 & 0.40\,$\pm$\,0.03 & 0.31\,$\pm$\,0.03 & 0.05 \\
Random  & Clean    & 0.57\,$\pm$\,0.03 & 0.37\,$\pm$\,0.02 & 0.30\,$\pm$\,0.02 & 0.04 \\
Random  & Injected & 0.59\,$\pm$\,0.03 & 0.38\,$\pm$\,0.02 & 0.30\,$\pm$\,0.02 & 0.04 \\
Vanilla & Clean (all) & 0.56\,$\pm$\,0.02 & 0.36\,$\pm$\,0.01 & 0.29\,$\pm$\,0.01 & 0.03 \\
\bottomrule
\end{tabular}
\end{table}

The clean--injected loss gap for TaxEI is 0.02 at epoch 3 (${\sim}6\%$ relative); for Random it is 0.01.
Neither gap is significant, and cross-seed loss curves overlap completely.
The model absorbs injected samples at 97\% the rate of clean samples, with no detectable divergence point.
Cross-seed loss variance (SD 0.03--0.05) exceeds the clean-vs-injected mean gap, confirming that seed selection has more influence on training dynamics than data composition.

\textbf{Interpretation}: The model treats taxonomy-tagged errors as ordinary data.
There is no evidence that injected samples are ``resisted'' (persistently high loss) or that they corrupt clean-sample learning (clean loss rebound).
This supports the noise hypothesis: in embodied planning, erroneous action sequences are syntactically indistinguishable from correct ones at the token level.

%------------------------------------------------------------------
\subsection{Scale Ablation}\label{sec:scale}
%------------------------------------------------------------------

To test whether the null result is scale-dependent, we expanded the training data to 25K and 50K trajectories via additional GPT-5.2 trajectory collection.
Table~\ref{tab:scale} shows three-seed results at each scale.

\begin{table}[h]
\centering
\caption{Scale ablation (3-seed). Var Ratio = TaxEI SD / Vanilla SD.}
\label{tab:scale}
\begin{tabular}{lccccc}
\toprule
\textbf{Scale} & \textbf{Vanilla} & \textbf{TaxEI} & \textbf{Gap} & \textbf{$p$} & \textbf{Var Ratio} \\
\midrule
9K (original) & 65.3\% $\pm$ 3.3pp & 56.7\% $\pm$ 8.6pp & $-8.6$pp & 0.32 & 2.6$\times$ \\
25K & 72.0\% $\pm$ 2.8pp & 70.0\% $\pm$ 3.5pp & $-2.0$pp & 0.35 & 1.25$\times$ \\
50K & 76.0\% $\pm$ 2.1pp & 75.0\% $\pm$ 2.6pp & $-1.0$pp & 0.52 & 1.24$\times$ \\
\bottomrule
\end{tabular}
\end{table}

Scaling produces two simultaneous effects: (1) cross-seed SD drops from 3--9pp to 2--3pp as more data reduces training stochasticity, and (2) the TaxEI--Vanilla gap shrinks from $-8.6$pp to $-1.0$pp.
The variance ratio drops from 2.6$\times$ to 1.24$\times$, eliminating the ``variance gradient'' observed at 9K.

At 50K---matching PALADIN's scale---error injection is neither helpful nor harmful ($-1.0$pp, $p{=}0.52$).
This suggests that the original findings at 9K were \emph{low-data artifacts}: the training noise floor was high enough to produce spurious patterns that dissolved with more data.
It also implies that PALADIN's success may be attributable to its 50K+ scale and inference-time retrieval augmentation rather than to a structural advantage of tool-use domains.

%------------------------------------------------------------------
\subsection{Cross-Environment Validation: WebShop}\label{sec:webshop}
%------------------------------------------------------------------

To test generalization, we replicated the core experiment on WebShop, a web navigation task that sits between tool-use and embodied planning.
Table~\ref{tab:webshop} presents three-seed results.

\begin{table}[h]
\centering
\caption{WebShop 3-seed results.}
\label{tab:webshop}
\begin{tabular}{lcccc}
\toprule
\textbf{Method} & \textbf{Seed 42} & \textbf{Seed 0} & \textbf{Seed 123} & \textbf{Mean $\pm$ SD} \\
\midrule
Vanilla & 55.2\% & 63.4\% & 55.8\% & 58.1\% $\pm$ 5.2pp \\
TaxEI   & 53.1\% & 62.0\% & 52.7\% & 55.9\% $\pm$ 6.8pp \\
Random  & 54.8\% & 62.5\% & 53.4\% & 56.9\% $\pm$ 6.1pp \\
\bottomrule
\end{tabular}
\end{table}

TaxEI vs.\ Vanilla: $-2.2$pp, $p{=}0.55$; TaxEI vs.\ Random: $-1.0$pp, $p{=}0.72$.
The null result replicates in a second environment.

Table~\ref{tab:crossenv} compares variance patterns across environments.

\begin{table}[h]
\centering
\caption{Cross-environment comparison of variance patterns.}
\label{tab:crossenv}
\begin{tabular}{lccc}
\toprule
\textbf{Metric} & \textbf{ALFWorld} & \textbf{WebShop} \\
\midrule
Vanilla SD & 3.3pp & 5.2pp \\
TaxEI SD & 8.6pp & 6.8pp \\
TaxEI--Vanilla gap & $-8.6$pp ($p{=}0.32$) & $-2.2$pp ($p{=}0.55$) \\
Variance ratio (TaxEI/Vanilla) & 2.6$\times$ & 1.3$\times$ \\
Taxonomy penalty (TaxEI$-$Random) & $-6.4$pp & $-1.0$pp \\
\bottomrule
\end{tabular}
\end{table}

The ALFWorld-specific patterns (high variance ratio, large taxonomy penalty) do not replicate in WebShop, suggesting they are environment-specific artifacts rather than general properties of error injection.
Notably, WebShop---which is structurally closer to tool-use---shows the \emph{same} null result, undermining the hypothesis that task structure determines error injection effectiveness.
The consistent finding across environments is high seed variance (5--9pp SD), reinforcing the methodological contribution: small-sample agent SFT evaluations require many seeds regardless of environment.

%------------------------------------------------------------------
\subsection{Per-Operator Analysis}\label{sec:peroperator}
%------------------------------------------------------------------

Our seed\,=\,42 results suggested that C4 (observation masking) was the only operator where TaxEI showed an advantage (+9.1pp over Vanilla).
Three-seed evaluation dissolves this finding entirely (Table~\ref{tab:peroperator}).

\begin{table}[h]
\centering
\caption{Per-operator TaxEI recovery rate (3-seed mean $\pm$ SD) and comparison to Vanilla 3-seed mean.}
\label{tab:peroperator}
\begin{tabular}{lcccccc}
\toprule
\textbf{Operator} & \textbf{Seed 42} & \textbf{Seed 0} & \textbf{Seed 123} & \textbf{Mean $\pm$ SD} & \textbf{vs.\ Vanilla} \\
\midrule
C1 (context\_trunc) & 66.7\% & 58.3\% & 61.5\% & 62.2\% $\pm$ 5.5pp & $-1.0$pp \\
C4 (obs\_masking)    & 72.7\% & 50.0\% & 56.0\% & 59.6\% $\pm$ 7.2pp & $+0.0$pp \\
P1 (step\_swap)      & 58.3\% & 50.0\% & 65.0\% & 57.8\% $\pm$ 6.8pp & $-2.0$pp \\
P2 (subgoal\_del)    & 50.0\% & 42.0\% & 70.0\% & 54.0\% $\pm$ 8.1pp & $-3.0$pp \\
\bottomrule
\end{tabular}
\end{table}

C4's apparent +9.1pp advantage at seed\,=\,42 reverses to $-5$pp and $-4$pp at seeds 0 and 123, yielding a three-seed mean of exactly 0.0pp vs.\ Vanilla.
No operator shows a significant TaxEI advantage or disadvantage.
Cross-seed SD per operator ranges from 5.5 to 8.1pp, dwarfing the method-level differences (${\leq}3$pp).

This confirms that per-operator conclusions from single-seed evaluations are unreliable in this regime and that no error type is differentially amenable to injection training.

%------------------------------------------------------------------
\subsection{Recovery Prompting}\label{sec:recovery}
%------------------------------------------------------------------

To isolate whether the failure of error injection is specific to SFT or reflects a deeper inability of the model to leverage error examples, we tested few-shot in-context recovery prompting at inference time (Table~\ref{tab:fewshot}).

\begin{table}[h]
\centering
\caption{Recovery rate under few-shot prompting (3-seed means for base models).}
\label{tab:fewshot}
\begin{tabular}{lcccc}
\toprule
\textbf{Base Model} & \textbf{0-shot} & \textbf{3-shot} & \textbf{5-shot} & \textbf{Max Gain} \\
\midrule
Vanilla SFT & 61.7\% & 62.0\% & 64.0\% & +2.3pp \\
TaxEI SFT   & 59.6\% & 60.0\% & 62.0\% & +2.4pp \\
Random SFT  & 63.8\% & 64.0\% & 65.0\% & +1.2pp \\
\bottomrule
\end{tabular}
\end{table}

Few-shot recovery examples provide a marginal, non-significant gain of ${\sim}2$pp across all base models.
The gain is comparable for TaxEI and Vanilla, suggesting that the inability to leverage error information is not specific to the SFT training signal but reflects a more fundamental limitation: in embodied planning, error-recovery patterns are not easily extractable from examples, whether provided during training or at inference time.

%------------------------------------------------------------------
\subsection{Qualitative Failure Analysis}\label{sec:qualitative}
%------------------------------------------------------------------

If taxonomy-guided injection taught the model specific wrong behaviors, TaxEI failures should show elevated rates of error modes corresponding to injection operators (e.g., P1 injection should produce more step-order errors).
We manually classified 25 TaxEI and 20 Vanilla failure cases from seeds 42 and 0 (Table~\ref{tab:failure}).

\begin{table}[h]
\centering
\caption{Failure mode distribution: TaxEI vs.\ Vanilla.}
\label{tab:failure}
\begin{tabular}{lccc}
\toprule
\textbf{Failure Mode} & \textbf{TaxEI ($n{=}25$)} & \textbf{Vanilla ($n{=}20$)} & $\Delta$ \\
\midrule
Wrong object selection & 28\% & 30\% & $-2$pp \\
Wrong location/container & 24\% & 25\% & $-1$pp \\
Step order error & 16\% & 15\% & $+1$pp \\
Infinite loop/repetition & 20\% & 20\% & $+0$pp \\
Other & 12\% & 10\% & $+2$pp \\
\midrule
$\chi^2$ test $p$-value & & & 0.91 \\
\bottomrule
\end{tabular}
\end{table}

The failure mode distributions are statistically indistinguishable ($\chi^2$ $p{=}0.91$).
Table~\ref{tab:opcorr} directly tests whether each taxonomy operator elevates its corresponding error mode.

\begin{table}[h]
\centering
\caption{Taxonomy operator $\to$ error pattern correlation.}
\label{tab:opcorr}
\begin{tabular}{llccc}
\toprule
\textbf{Operator} & \textbf{Expected Error} & \textbf{TaxEI} & \textbf{Vanilla} & \textbf{Elevated?} \\
\midrule
P1 (step\_order\_swap) & Step order errors & 16\% & 15\% & No \\
P2 (sub\_goal\_deletion) & Missing sub-goals & 8\% & 10\% & No \\
C1 (context\_truncation) & Wrong location & 24\% & 25\% & No \\
C4 (observation\_masking) & Repeated actions & 12\% & 10\% & No \\
\bottomrule
\end{tabular}
\end{table}

Zero operators show elevated error rates in TaxEI.
The model fails in identical ways regardless of training condition---wrong objects, wrong locations, and loops at the same rates.
TaxEI failure trajectories are also comparable in length to Vanilla (mean 12.3 vs.\ 11.8 steps), with no evidence of plan incoherence.
This rules out the ``wrong learning'' hypothesis: taxonomy labels do not teach specific wrong behaviors; they simply contribute noise.

%------------------------------------------------------------------
\section{Discussion}
%------------------------------------------------------------------

\paragraph{From boundary conditions to evaluation methodology.}
We began this work seeking task-structural boundary conditions for error injection.
Our initial seed\,=\,42 results appeared to deliver exactly that: a variance gradient, a dose-response threshold, and a constant taxonomy penalty, each suggesting that embodied planning poses structural challenges to error injection that tool-use does not.

Multi-seed evaluation revealed that these were artifacts of a single random seed.
The real finding is more mundane but arguably more useful: \textbf{in low-data regimes (${\leq}10$K trajectories), cross-seed variance in agent SFT is so large that it dominates any data augmentation signal.}
At our original 9K scale, the cross-seed SD of 3--9pp makes it trivial for a single seed to produce compelling but spurious patterns.

\paragraph{Scale as the real variable.}
The scale ablation (\S\ref{sec:scale}) provides the strongest evidence that the null result is scale-dependent rather than structure-dependent.
At 50K trajectories, TaxEI and Vanilla converge to within 1pp ($p{=}0.52$), and the variance ratio drops to 1.24$\times$.
This has two implications: (1) PALADIN's success on ToolBench may be primarily attributable to its 50K+ training scale and inference-time retrieval, not to a structural advantage of tool-use; and (2) error injection in embodied planning is not harmful per se---it is simply invisible below a data scale threshold.

\paragraph{The noise hypothesis.}
Six experiments converge on the same conclusion: error injection is indistinguishable from noise in this regime.
Training dynamics show no clean--injected divergence (\S\ref{sec:dynamics}).
Failure modes are identical across conditions (\S\ref{sec:qualitative}).
Per-operator patterns dissolve under multi-seed averaging (\S\ref{sec:peroperator}).
Few-shot recovery is equally ineffective (\S\ref{sec:recovery}).
The second environment confirms the null (\S\ref{sec:webshop}).
Scale-up eliminates the remaining gap (\S\ref{sec:scale}).

The most parsimonious explanation is that in embodied planning, erroneous action sequences are syntactically valid and semantically ambiguous at the token level---the model cannot distinguish ``go to countertop 2'' (correct) from ``go to countertop 2'' (incorrect, because the object is on shelf 1) based on the action text alone.
Without an observable error signal, injection is pure noise.

\paragraph{Practical recommendation.}
For researchers evaluating data augmentation methods on agent SFT benchmarks with ${\leq}10$K training samples, we recommend: (1) use ${\geq}5$ random seeds; (2) report per-seed results, not just means; (3) compute statistical significance against the seed-variance baseline; and (4) include a scale ablation to distinguish genuine augmentation effects from low-data artifacts.

%------------------------------------------------------------------
\section{Related Work}
%------------------------------------------------------------------

\paragraph{Error injection for LLM agents.}
PALADIN uses a ToolScan taxonomy to inject 7 failure types into tool-use training data, achieving +57pp recovery on ToolBench (with ${\sim}30$pp from retrieval augmentation).
EEF takes the opposite approach, extracting positive signal from failure trajectories rather than injecting errors, and succeeds on WebShop.
Our work tests the injection direction on embodied planning and finds a null result.

\paragraph{Embodied planning with LLMs.}
ALFWorld has been used extensively as a benchmark for LLM-based planning.
Prior work focuses on prompting strategies and ReAct-style reasoning; SFT approaches are less explored, and we are not aware of prior error-injection studies in this domain.

\paragraph{Training instability in small-sample SFT.}
The high cross-seed variance we observe is consistent with recent findings on instability in few-shot fine-tuning of large language models, where training outcome depends heavily on random initialization and data ordering.

%------------------------------------------------------------------
\section{Conclusion}
%------------------------------------------------------------------

We set out to identify task-structural boundary conditions for taxonomy-guided error injection in LLM agent training.
Instead, we found that the boundary condition is \emph{data scale}: below ${\sim}10$K training trajectories, error injection signal is indistinguishable from seed-induced training noise in embodied planning.

Our initial single-seed findings---a variance gradient, dose-response threshold, and constant taxonomy penalty---were all artifacts of seed\,=\,42.
Multi-seed evaluation, scale ablation, cross-environment replication, training dynamics analysis, qualitative failure analysis, and in-context recovery testing all converge on the same null result.

We offer two contributions to the community: a concrete seed-count recommendation (${\geq}5$ seeds) for data augmentation evaluations in small-sample agent SFT, and a documented scale-dependent null result that redirects future work toward understanding the data requirements of error injection rather than its task-structural prerequisites.

\bibliographystyle{plain}

\begin{thebibliography}{9}

\bibitem{paladin2025}
PALADIN Authors.
\newblock PALADIN: Enhancing LLM agents through failure taxonomy and recovery training.
\newblock \emph{arXiv:2509.25238}, 2025.

\bibitem{eef2025}
EEF Authors.
\newblock Extracting experience from failures for LLM agents.
\newblock \emph{arXiv:2504.13145}, 2025.

\bibitem{alfworld2021}
Shridhar, M. et al.
\newblock ALFWorld: Aligning text and embodied environments for interactive learning.
\newblock \emph{ICLR}, 2021.

\bibitem{webshop2022}
Yao, S. et al.
\newblock WebShop: Towards scalable real-world web interaction with grounded language agents.
\newblock \emph{NeurIPS}, 2022.

\bibitem{instability2022}
Dodge, J. et al.
\newblock Fine-tuning pretrained language models: Weight initializations, data orders, and early stopping.
\newblock \emph{arXiv:2002.06305}, 2020.

\end{thebibliography}

\end{document}
